% sections/chapter-findings.tex
\chapter{Empirical Findings}
\label{ch:findings}

The empirical findings from the eighteen months of observation at Nexa, Synapse, and Vertex reveal a systematic divergence between formal coordination processes and actual attentional allocation. This chapter details the quantitative results of the Ambient Coordination Load Index (ACLI) and explores the qualitative mechanisms of what we term \emph{reconciliation drift}.

\section{ACLI Distribution across Sites}

The primary quantitative result is the distribution of the ACLI across the three observed organizations. As shown in Figure~\ref{fig:acli_dist}, the median ACLI varies significantly with organization size and tooling maturity.

\begin{figure}[htbp]
\centering
\begin{tikzpicture}[scale=0.85]
  % Draw axes
  \draw[thick, ->, color=inkblue] (0,0) -- (10,0) node[right, font=\small] {Weekly Mean ACLI};
  \draw[thick, ->, color=inkblue] (0,0) -- (0,5) node[above, font=\small] {Density};
  
  % Plot curves
  % Nexa (dashed, inkblue)
  \draw[thick, color=inkblue, dashed] plot [smooth, tension=0.8] coordinates {
    (0.5, 0.1) (2.0, 0.8) (3.5, 2.2) (5.0, 1.2) (7.0, 0.4) (9.0, 0.05)
  };
  \node[color=inkblue, right] at (5.2, 1.5) {\small Nexa};

  % Synapse (solid, claret)
  \draw[thick, color=claret] plot [smooth, tension=0.8] coordinates {
    (0.5, 0.05) (1.5, 0.4) (2.8, 1.1) (4.2, 2.8) (5.5, 1.8) (7.5, 0.5) (9.0, 0.02)
  };
  \node[color=claret, right] at (4.5, 2.9) {\small Synapse};

  % Vertex (dotted, slategray)
  \draw[thick, color=slategray, dotted, line width=1.2pt] plot [smooth, tension=0.8] coordinates {
    (0.5, 0.02) (2.5, 0.5) (4.5, 1.2) (6.0, 2.5) (7.5, 1.8) (9.0, 0.3) (9.5, 0.05)
  };
  \node[color=slategray, right] at (6.5, 2.6) {\small Vertex};

\end{tikzpicture}
\caption{Estimated density distributions of the Ambient Coordination Load Index (ACLI) across the three field sites.}
\label{fig:acli_dist}
\end{figure}

The distribution demonstrates that at Vertex, which has the largest headcount and the most fragmented tool stack, the average ACLI is shifted rightward, indicating that team members spend a significantly larger portion of their day resolving coordination discrepancies.

\section{The Dynamics of Reconciliation Drift}

Reconciliation drift\index{reconciliation drift} occurs when team members stop relying on the formal systems of record (e.g., ticket trackers, project wikis) because they believe the updates are stale. This creates a self-reinforcing cycle:

\begin{enumerate}
  \item A team member observes a discrepancy between a ticket state and reality.
  \item They initiate a chat thread to confirm the actual status.
  \item This confirms the ticket is stale, reducing trust in the system of record.
  \item Subsequent status checks bypass the system of record entirely, moving to chat.
  \item The ticket system falls further out of date, cementing the drift.
\end{enumerate}

At Synapse, where the team is small (21 members), this drift was managed through high-frequency chat communication. However, at Vertex (46 members), the volume of chat threads required to maintain a current map of ownership led to severe attentional fragmentation, with the median engineer checking chat channels 42 times per day.

\begin{fieldexcerpt}[title={Field note, Synapse, week 12}]
``We don't really use the project board anymore except to show the executives what we did last month. If I want to know what Elena is working on, I just check the daily standup thread or message her directly. It's faster, and it's actually accurate.''\\[2pt]
\hfill--- Senior Machine Learning Researcher
\end{fieldexcerpt}

The consequence of this drift is not merely that coordination work is moved to a different channel, but that it changes character. What was once a structural representation of the project becomes a set of transient, unstructured text streams that must be read and synthesized continuously to avoid missing critical updates. This shift from structural to conversational coordination is the primary driver of ambient coordination load.
